\section{A 201-variable natural quartic}

All coordinates in this section range independently over
$\mathbb N=\{0,1,2,\ldots\}$. The principal polynomial is the sum of the
squares of $153$ literal quadratic residuals. The defining residual list
and its complete collected polynomial are supplied respectively in
\begin{center}\small
\texttt{natural201/equation.json}\\
\texttt{natural201/expanded\_polynomial.json}
\end{center}
It has $201$ variables, exact
degree $4$, $6002$ nonzero monomials and coefficient height $1860103$.
The inherited fixed-calculus calibration is unchanged: it has a natural zero
if and only if the specified effectively presented ZFC calculus has a finite
syntactic refutation. Its no-natural-zero sentence thus retains the inherited
$\operatorname{Con}(\mathrm{ZFC})$ interpretation. This is not a consistency
proof, an optimality claim, or a new end-to-end kernel formalization.

\begin{center}
\begin{tabular}{lrrrrr}
Profile & Variables & Rows & Degree & Terms & Height\\\hline
Principal &201&153&4&6002&1860103\\
Affine parent &202&154&4&6000&1854571\\
Sparse sextic &165&117&6&6558&1860101\\
Smaller sextic &164&116&6&15774&1860101\\
Sparse octic &139&91&8&11873&1117699986\\
Smaller octic &138&90&8&16515&2933403776
\end{tabular}
\end{center}

\paragraph{Natural affine charts certified modulo source equations.}
Let $R=(R_0,\ldots,R_{m-1})$ be integral polynomials in $n$ coordinates.
Suppose a set $S$ of coordinates has retained source equations
\[
 R_{j(v)}=X_v-f_v(X),\qquad f_v\in\mathbb N[X],\quad v\in S,
\]
and an ordering in which every coordinate of $S$ appearing in $f_v$ is
earlier than $v$. Let $F$ be an integral affine automorphism, with integral
affine inverse $G$, fixing the coordinates outside $S$. Suppose that for each
$v\in S$ a certificate supplies
\[
 F_v=P_v+\sum_j a_{vj}R_j,\qquad P_v\in\mathbb N[X],\quad a_{vj}\in\mathbb Z.
 \tag{N1}
\]
Then $F$ and $G$ give inverse bijections between the natural zero sets of
$R$ and $R\circ G$.

Indeed, at a natural zero $x$ of $R$, equation~(N1) proves $F_v(x)\ge0$;
the other coordinates are unchanged. Conversely, if $y$ is a natural zero
of $R\circ G$, put $z=G(y)$ in $\mathbb Z^n$. Coordinates outside $S$ are
already natural. The retained equations give $z_v=f_v(z)$, so induction in
the stated source order proves every remaining coordinate nonnegative.
The two affine inverse identities prove the inverse laws. Crucially, no
claim is made that $G$ preserves the entire natural orthant.
An elementary integral row operation $R_i\leftarrow R_i+kR_j$, $i\ne j$,
is reversed by subtraction and does not change the common zero set.
Affine substitutions do not increase degree; applying the affine inverse
proves preservation of maximum residual degree.

\paragraph{The literal chart.}
The source is the preceding pooled $202$-coordinate natural system, with
its full source closure and proof retained in \texttt{inherited/}.
A positive syntax-index-budget accumulator, the twelve earlier syntax
shears, and forty further literal affine shears yield the selected chart;
twenty-four recorded elementary row shears are then applied. Complete replay
inputs and both affine maps occur under \texttt{chart202/}.
There are $33$ changed coordinates. Twenty-nine final forward maps are
coefficientwise nonnegative. The other four require only source rows
$13,17,34,75$, for coordinates $27,31,49,93$ respectively. For example,
\[
 F_{31}=X_{27}+X_{31}-16X_{25},\qquad F_{93}=X_{93}-X_{92}.
\]
Subtracting source row $17$ from $F_{31}$ gives the recorded polynomial
with nonnegative coefficients. Source row $75$ has the form
$X_{93}-HX_{92}$, where $H$ is coefficientwise nonnegative with constant
$236$; thus $F_{93}=(H-1)X_{92}+R_{75}$ is nonnegative on source zeros.
All $33$ original positive defining equations and their dependency order
are checked from literal polynomials. The inherited Pell and power
interpretations are used only after the original natural zero has been
restored; no lost orientation is assumed away.

\paragraph{The new quadratic contraction.}
In the final $202$-coordinate chart, before renumbering, the only rows
containing $X_{57}$ are
\[
 R_{42}=X_{57}-2-6X_{48}-4X_{48}^{2},\qquad
 R_{44}=X_{59}-4X_{57}-X_{58}-X_{58}^{2}.
\]
The coordinated shear of $X_{59}$ has removed the other incidences.
Substitute $X_{57}=2+6X_{48}+4X_{48}^{2}$, and delete coordinate $57$
and row $42$. The remaining occurrence becomes
\[
 X_{59}-8-24X_{48}-16X_{48}^{2}-X_{58}-X_{58}^{2}.
\]
Every residual remains quadratic. The restoring polynomial is nonnegative
on every natural retained assignment, and the deleted monic row forces
precisely that value. Deletion and extension are therefore inverse natural
zero-set maps. The resulting sum of $153$ squares is the asserted
$201$-variable quartic. The literal collected coefficients establish exact
degree, support and height. No lower-bound theorem is inferred from the
finite search for further contractions in this particular row basis.

\paragraph{Signed monic contractions for the degree tradeoffs.}
There is a more general certificate which avoids requiring every restoring
polynomial to have nonnegative coefficients. Start with a chart
$C=U(R\circ G)$, where $U$ is an invertible integral residual matrix.
Repeatedly select a row $aY_v+h(Y_{\ne v})$, $a\in\{1,-1\}$, substitute
$Y_v=-ah$ into the remaining rows and delete that row and coordinate.
These operations give a polynomial extension $E$ of the retained coordinates
and a projection $\pi$, inverse to one another on the integer zero sets.
This follows inductively because each deleted monic equation forces its
restoring value, without division or an additional congruence.

Let $D$ be the set of deleted source-coordinate indices. Assume the original
system has positive acyclic defining equations for every coordinate in
$S\cup D$. Given a natural retained zero $y$, extend to $E(y)$ over the
integers and put $z=G(E(y))$. All original residuals vanish. Coordinates
outside $S\cup D$ remain unchanged natural coordinates. Induction along
the original positive definitions proves $z\in\mathbb N^n$. Equation~(N1)
then gives $F(z)=E(y)\in\mathbb N^n$. The inverse direction is projection
of the natural chart image of an original natural zero. Hence the same
maps are inverse bijections on natural zero sets, even when intermediate
restoring formulas have negative coefficients.

The $165$-, $164$-, and $139$-variable companions contract $37$, $38$, and
$63$ coordinates from the new chart; their certified restoration sets have
$69$, $70$, and $83$ coordinates. The $138$-variable companion contracts
$64$ coordinates from the preceding-round chart and has a $74$-coordinate
restoration set. In the degree-eight profiles, the deleted original outer
$x$ is restored from its explicit positive equation
$x=1+\mathrm{endpoint.ut}(1+\mathrm{mul.8})$; no sign of a free integer
is postulated. The equivalence is existential, not a new assertion with
$x$ held as an externally fixed parameter. Every contraction equation,
substitution, restoration order and retained-index renumbering is supplied.
The residual degrees are at most three or four, giving the displayed
sextic and octic tradeoffs. They are not quartic-dimension improvements.

\paragraph{Current-source comparison and verification boundary.}
At the pinned current snapshot
\texttt{ecb9a4b4dda617512952525d50f17343700aac52}, the authenticated
source178 principal and sparse natural $203$-variable quartics have
$155$ rows and respectively $8237$ terms/height $799349232$ and
$8263$ terms/height $207116893$. The new $201$-variable quartic improves
all four measures against both. The source178 low-height alternative
has $49682$ terms and height $481920$; it remains a height tradeoff.
Literal current-source equations and a full independent recollection are
included, so this is not a comparison with an obsolete summary.

The preceding principal had $202$ variables, $154$ rows and $9353$ terms.
The new counts are $201$, $153$ and $6002$. Height rises from $1854167$ to
$1860103$, an increase of $5936$ (approximately $0.32\%$). Accordingly,
we do not claim simultaneous improvement of every previous metric.
The inherited low-height $202$-variable form, with $16927$ terms and
height $416488$, is preserved. The new $165$-variable sextic strictly
improves every displayed metric over the earlier $200$-variable octic, which
had $152$ rows, $10392$ terms and height $13654244$.

The standard-library-only builder reconstructs all literal rows and
coefficients. An independent exponent-count implementation verifies the
affine inverse identities, positive representatives, original restoration
DAGs, monic eliminations, every residual and every expanded coefficient.
Negative controls reject corrupted representatives, gate orders, eliminated
definitions and coefficients. Normal and optimized Python execution,
byte-identical temporary reconstruction, LaTeX round-trips, current-source
recollection and replay of the complete inherited finite construction closure
are recorded. These are finite exact certificates accompanying the ordinary
unbounded semantic and logical proofs, not a claim of a new kernel proof.
